% ----- CHAUPAI 30 -----

\newpage

\subsection*{\deva{त्रिंशत्तम चौपाई} (Thirtieth Chaupai)}
\addcontentsline{toc}{subsection}{\deva{त्रिंशत्तम चौपाई} (Thirtieth Chaupai) — \deva{साधु सन्त...}}

\begin{minipage}[t]{0.55\textwidth}
\vspace{0pt}

{\large \linespread{1.5}\selectfont
\deva{साधु सन्त के तुम रखवारे।}\\
\deva{असुर निकंदन राम दुलारे॥}
\par}

\vspace{0.5em}

{\large \linespread{1.5}\selectfont
\telu{సాధు సన్త కే తుమ రఖవారే।}\\
\telu{అసుర నికందన రామ దులారే॥}
\par}

\vspace{0.5em}

{\large \linespread{1.3}\selectfont
\textit{sādhu santa ke tuma rakhavāre |}\\
\textit{asura nikandana rāma dulāre ||}
\par}
\end{minipage}%
\hfill
\begin{minipage}[t]{0.42\textwidth}
\vspace{0pt}
\begin{tcolorbox}[anchorbox]
\textbf{The Righteous Protector:} He is the gentle guardian of the innocent (\textit{Sadhu}) and the fierce uprooter of toxic elements (\textit{Asura Nikandana}). Students must learn to cultivate both immense compassion for others and the firm boundaries needed to stand up against negative influences.
\vspace{0.5em}\newline He struck the perfect balance: fiercely eliminating toxic threats while gently protecting the innocent in his care.\newline\null\hfill\href{https://www.valmiki.iitk.ac.in/sloka?field_kanda_tid=5&language=dv&field_sarga_value=46}{\textit{Sundara Kanda, 46:25-27}}
\end{tcolorbox}
\end{minipage}

\vspace{0.5em}
\subsubsection*{\textbf{Visual Rhythm Map:}}
\begin{table}[h!]
\centering
\renewcommand{\arraystretch}{1.2}
\setlength{\tabcolsep}{0pt}
\begin{tabularx}{\textwidth}{|*{16}{>{\centering\arraybackslash\hsize=1.0\hsize}X|}}
\hline
\multicolumn{3}{|>{\centering\arraybackslash\hsize=3.0\hsize}X|}{\textbf{\guru{सा}\laghu{धु}} \par (3)} &
\multicolumn{3}{>{\centering\arraybackslash\hsize=3.0\hsize}X|}{\textbf{\guru{स}\laghu{न्त}} \par (3)} &
\multicolumn{2}{>{\centering\arraybackslash\hsize=2.0\hsize}X|}{\textbf{\guru{के}} \par (2)} &
\multicolumn{2}{>{\centering\arraybackslash\hsize=2.0\hsize}X|}{\textbf{\laghu{तु}\laghu{म}} \par (2)} &
\multicolumn{6}{>{\centering\arraybackslash\hsize=6.0\hsize}X|}{\textbf{\laghu{र}\laghu{ख}\guru{वा}\guru{रे}} \par (6)} \\
\hline
\end{tabularx}
\vspace{-1.5em}
\end{table}

\begin{table}[h!]
\centering
\renewcommand{\arraystretch}{1.2}
\setlength{\tabcolsep}{0pt}
\begin{tabularx}{\textwidth}{|*{16}{>{\centering\arraybackslash\hsize=1.0\hsize}X|}}
\hline
\multicolumn{3}{|>{\centering\arraybackslash\hsize=3.0\hsize}X|}{\textbf{\laghu{अ}\laghu{सु}\laghu{र}} \par (3)} &
\multicolumn{5}{>{\centering\arraybackslash\hsize=5.0\hsize}X|}{\textbf{\laghu{नि}\guru{कं}\laghu{द}\laghu{न}} \par (5)} &
\multicolumn{3}{>{\centering\arraybackslash\hsize=3.0\hsize}X|}{\textbf{\guru{रा}\laghu{म}} \par (3)} &
\multicolumn{5}{>{\centering\arraybackslash\hsize=5.0\hsize}X|}{\textbf{\laghu{दु}\guru{ला}\guru{रे}} \par (5)} \\
\hline
\end{tabularx}
\vspace{-1.5em}
\end{table}

\vspace{1em}
\textbf{ \deva{सरल-संस्कृतम्}:}\\
 \deva{साधूनां सन्तानां च भवान् रक्षकः अस्ति।}\\
 \deva{भवान् असुराणां विनाशकः, श्रीरामस्य प्रियः च अस्ति॥}\\


You are the protector of sages and saints.\\
You are the destroyer of demons, and the darling of Lord Rama.


\vspace{1em}
\newpage
\textbf{\deva{प्रतिपदार्थः} (Meaning of each word):}
\begin{table}[h!]
\centering
\renewcommand{\arraystretch}{1.2}
\begin{tabularx}{\textwidth}{@{} l l X X @{}}
\toprule
\textbf{Awadhi} & \textbf{Sanskrit} & \textbf{English} & \textbf{Grammar \& Notes} \\
\midrule
\deva{साधु सन्त के} / \telu{సాధు సన్త కే}  &  \deva{साधु-सन्तानाम्}  &  Of sages and saints  &   \\
\deva{तुम} / \telu{తుమ}  &  \deva{त्वम्}  &  You  &   \\
\deva{रखवारे} / \telu{రఖవారే}  &  \deva{रक्षकः}  &  Protector  & \\ 
\deva{असुर} / \telu{అసుర}  &  \deva{असुराणाम्}  &  Of demons  &   \\
\deva{निकंदन} / \telu{నికందన}  &  \deva{विनाशकः / निकृन्दन}  &  Destroyer / Uprooter  &   \\
\deva{राम दुलारे} / \telu{రామ దులారే}  &  \deva{राम-प्रियः}  &  Dear to Rama / Darling  &   \\
\bottomrule
\end{tabularx}
\end{table}

\newpage
